%----------------------------------------------------------------------
% 緒論
%----------------------------------------------------------------------

\chapter{Introduction}\label{chap:intro}

In recent years, to mitigate the inference latency caused by token-by-token generation in autoregressive large language models, researchers have begun exploring approaches that predict multiple future tokens within a single model inference step. Gloeckle et al.\cite{gloeckle2024mtp} introduced Multi-Token Prediction (MTP), which enables the model to simultaneously predict tokens at multiple future positions during pre-training through a shared model trunk and multiple prediction heads. Subsequently, Samragh et al.\cite{samragh2025future} further investigated how multi-token prediction capabilities can be introduced into existing pretrained language models through supervised fine-tuning and additional consistency objectives. In addition, DeepSeek-V3\cite{deepseekai2024deepseekv3} incorporated MTP into the pre-training objective of a large-scale language model, demonstrating the practical potential of multi-token prediction for improving inference efficiency.

\section{Background and Motivation}\label{sec:1-motivation}

\subsection{Background}

Large language models typically generate text in an autoregressive manner. Given an input sequence, the model predicts the next token based on the previously observed tokens, appends the newly generated token to the sequence, and then performs the next prediction step. Although this next-token prediction paradigm preserves a clear causal dependency, inference generally requires tokens to be generated sequentially. As a result, the number of forward passes increases with the length of the output sequence.

Multi-Token Prediction (MTP) extends conventional next-token prediction to multiple future positions. An MTP model consists of a shared model trunk and multiple prediction heads. Given the same input representation, the first prediction head predicts the next token, while the remaining prediction heads are responsible for predicting tokens at more distant future positions. Consequently, multiple future-token predictions can be obtained from a single forward pass through the shared model trunk. Since most of the computational cost is concentrated in the shared trunk, while the additional cost introduced by multiple prediction heads is relatively limited, MTP provides the potential to reduce redundant model computation and accelerate decoding.

During inference, the future tokens generated by multiple prediction heads can be treated as draft tokens. If these draft tokens are consistent with the subsequent verification results, multiple tokens can be accepted within a single decoding step, thereby reducing the number of model forward passes required by conventional token-by-token generation. Therefore, the accuracy of future-token prediction and the reliability of the outputs produced by different prediction heads directly affect the potential decoding efficiency of MTP.

\subsection{Motivation}

Although Multi-Token Prediction (MTP) provides a promising approach for accelerating autoregressive decoding, effectively optimizing multiple prediction heads remains challenging. In standard MTP training, each prediction head is independently optimized using the cross-entropy loss associated with its corresponding ground-truth token. While this objective allows each head to learn future-token prediction, it does not explicitly enforce consistency among the predictions produced by different heads.

As the prediction horizon increases, the uncertainty of future-token prediction also increases. Consequently, prediction heads responsible for more distant future positions are generally more difficult to optimize and tend to exhibit lower prediction accuracy than the first prediction head. Since the effectiveness of multi-token decoding depends on the quality and reliability of the predicted draft tokens, inaccurate predictions from later heads may reduce the number of draft tokens that can be successfully accepted during decoding.

Based on these observations, this study investigates whether additional consistency constraints can improve the optimization of future-token prediction heads without modifying the underlying MTP architecture. Specifically, the first prediction head, which performs next-token prediction, is used as a relatively reliable reference to guide the remaining prediction heads. By encouraging consistency across prediction heads at both the output probability distribution and latent representation levels, this study aims to improve the prediction quality of distant future-token heads and enhance the effectiveness of multi-token decoding.

\section{Contributions}\label{sec:1-contribution}

\begin{enumerate}
    \item We propose Selective Distribution Consistency for MTP training, which aligns the output distributions of student heads with the first prediction head while filtering low-confidence teacher predictions.

    \item We combine SDC with Latent Consistency Matching (LCM) to improve consistency across prediction heads at both the output distribution and latent representation levels without modifying the MTP architecture.

    \item Experiments on a 1.3B-parameter model show that SDC+LCM reduces Oracle PPL across prediction heads, increases Average Draft Tokens Accepted, and substantially reduces Generative PPL.

\end{enumerate}


% \section{章節概述\small{（列點的方式）}}

% \begin{itemize}
%     \item 在這種困難的抉擇下，本人思來想去，寢食難安。 要想清楚，章節概述，到底是一種怎麽樣的存在。 塞涅卡在不經意間這樣說過，生命如同寓言，其價值不在與長短，而在與內容。這啟發了我。
%     \item 那麽， 既然如此， 斯賓諾莎在不經意間這樣說過，最大的驕傲於最大的自卑都表示心靈的最軟弱無力。這不禁令我深思。
%     \item 要想清楚，章節概述，到底是一種怎麽樣的存在。 了解清楚章節概述到底是一種怎麽樣的存在，是解決一切問題的關鍵。
%     \item 就我個人來說，章節概述對我的意義，不能不說非常重大。 馮學峰曾經說過，當一個人用工作去迎接光明，光明很快就會來照耀著他。帶著這句話，我們還要更加慎重的審視這個問題。
% \end{itemize}